\section{Related Work}

We connect to four threads. (i) \textbf{EEG domain generalization / cross-subject MI}, especially the MOABB
benchmarking protocol~\citep{jayaram2018moabb}, from which our datasets and leave-one-subject-out (LOSO)
evaluation derive. (ii) \textbf{Graph EEG networks} --- DGCNN~\citep{song2018dgcnn},
RGNN~\citep{zhong2019rgnn}, and LGGNet~\citep{ding2021lggnet} --- which motivate graph/node/edge objects.
(iii) \textbf{Known-good convolutional MI decoders} --- EEGNet~\citep{lawhern2018eegnet} and
ShallowConvNet/DeepConvNet~\citep{schirrmeister2017deep} --- which we use only as sanity references for task
learnability. (iv) \textbf{Conditional-invariance / information-penalty} methods --- conditional
domain-invariant representations~\citep{li2018conditional}, classifier-based CMI / posterior-KL leakage
proxies~\citep{mukherjee2020ccmi}, and domain-adversarial or marginal-invariance baselines.

CIGL differs by (a) \emph{auditing} graph and node leakage explicitly with a permutation-null proxy before
regularizing, (b) operating strictly source-only with target labels evaluation-only, and (c) restricting to
a graph/node penalty on a task-capable static-adjacency backbone rather than claiming a general
edge/dynamic-graph method. A positioning grid and per-citation verification status are kept in
\texttt{../cigl/RELATED\_WORK\_MATRIX.md} and \texttt{../cigl/CITATION\_TODO\_QUEUE.md}.

% CMI-Trace P0.5 related-work correction (exact prose from paper/cmi_trace/MANUSCRIPT_INTEGRATION.md).
% arXiv IDs are cited inline as text (no \citep key added, because REFERENCES_DRAFT.bib lives outside
% paper/cigl_latex/ and this revision is scoped to that directory only).
{\sloppy
\paragraph{Concurrent subject-identity audits.}
Two concurrent audits of subject identity in EEG foundation models are complementary to ours but ask
different questions. \textbf{Lin, Wu, and Jung (arXiv:2606.06647)} study \emph{marginal} subject dominance
in frozen foundation-model features --- LEACE-based erasure, the spectral carrier of subject identity, and
the consistency of the task axis --- i.e.\ a geometry/erasure audit of the marginal subject signal.
\textbf{Tai (arXiv:2606.09189)} is a \emph{privacy / attribute-leakage} study auditing cross-encoder
attribute transfer (whether attributes leak across independently trained encoders), not a marginal
subject-variance audit. We therefore do \textbf{not} describe both as ``marginal subject-variance audits.''
Crucially, FMScope's cohort-conditioned \emph{frozen-feature} diagnostic (an oracle that fits the erasure
geometry on the whole cohort, including the held-out subject's features) and our strict source-only
\emph{unseen-subject deployment} ask different questions and \textbf{can both be correct}: cohort-conditioned
geometry can be erasable while the same axis is not identifiable, or removable, from source subjects alone
(P0.5 bridge, oracle vs source-only).\par}
